\subsection{Sharper intersections of the protected cycles}

\begin{proposition}[Proved: a proper overlap lies on one boundary edge]
\label{C:protected-nonadjacent}
Use the protected objects $P_+,P_-$, their marked circles $S_+,S_-$,
and their cycle generators from Theorem~\ref{C:protected-cycles}.
For every object $X$ of the repaired construction, every protected
arrow $h:P_\epsilon\to X$, and every arrow $g:Y\to X$, the following
stronger routing assertion holds. If $h(S_\epsilon)\cap g(Y)$ contains
the images of two vertices which are not adjacent in $S_\epsilon$,
then $h$ factors through $g$. Consequently, if $h$ does not factor
through $g$, that intersection is empty, a single vertex, or two
adjacent vertices. The boundary graphs of two distinct protected cells
in $G_X$ meet in an empty set, one vertex, or one edge.
\end{proposition}

\begin{proof}
We prove the routing assertion by stage induction. At the first stage
there is no arrow between the two distinct protected objects and no
nonidentity endomorphism of either. A proper incoming $V$-arrow meets
the marked circle in at most one vertex; a $W$-arrow meets it in the
endpoints of its one boundary edge. Thus the assertion holds at the
birth stage.

Suppose first that $X$ is a new quasi-filter target, with neighboring
maps $\alpha_i:T\to X$ and old pair $\phi_0,\phi_1:S\rightrightarrows
T$. The case $g=1_X$ is immediate. Otherwise cyclically rename the
copies so that $g=\alpha_0g^*$ and write $h=\alpha_jh^*$. Choose two
nonadjacent marked vertices $p,q\in S_\epsilon$ whose images lie in
$g(Y)$. If $j\ne0$, take the next index $j'$ on the shortest arc from
$j$ to $0$. Its total length is at most $2k\leq s-2n$, so (10) puts
$h(p),h(q)$ in the overlap of copies $j,j'$. By (11) that overlap is
$\alpha_j\phi_r(S)=\alpha_{j'}\phi_t(S)$ for appropriate
$r,t\in\{0,1\}$. Injectivity of $\alpha_j$ puts $h^*(p),h^*(q)$ in
$\phi_r(S)$. The old induction hypothesis, applied to the same two
nonadjacent vertices in the fixed source $P_\epsilon$, gives an old
arrow $l:P_\epsilon\to S$ with $h^*=\phi_r l$. Hence
$h=\alpha_{j'}\phi_t l$. The entire protected arrow has moved one
step towards copy zero. Repeating at most $2k$ times writes
$h=\alpha_0\widetilde h$. Injectivity then puts
$\widetilde h(p),\widetilde h(q)$ in $g^*(Y)$, and the old assertion
factors $\widetilde h$ through $g^*$. This factors $h$ through $g$.

At a disjoint-sum target, $g=1_X$ is again immediate. Otherwise any
nonempty overlap puts $h$ and $g$ in the same summand, where the old
assertion applies. This completes the induction for every stage and
every numerical choice. A subset of a cycle of length at least four
with no nonadjacent pair has at most two vertices, and a two-vertex
such subset is an edge's pair of endpoints.

Finally suppose two distinct protected arrows have two common marked
vertices. If these were nonadjacent in either source circle, the
assertion just proved would factor that protected arrow through the
other. There are no arrows between $P_+$ and $P_-$, and their only
endomorphisms are identities; such a factorization would therefore
make the arrows equal. Thus the common vertices must be adjacent in
both circles. Their two boundary $W$-arrows share two image points, so
the original AKR routing invariant and $\operatorname{End}(W)=\{1\}$
make those $W$-arrows equal. A further common marked vertex would
give a nonadjacent pair. The graph intersection is therefore exactly
one edge in this case, proving the final assertion.
\end{proof}

\begin{proposition}[Proved: protected chains supported on two arrow images]
\label{C:protected-two-images}
Let $P(X)\subseteq K(X)\subseteq C_1(G_X;\mathbb Z)$ be the subgroup
freely generated by the protected cycles in
Theorem~\ref{C:protected-cycles}. For any two arrows
$f:Y\to X$ and $g:Z\to X$, let $U\subseteq G_X$ be the union of
their image subgraphs. Then
\[
 P(X)\cap C_1(U;\mathbb Z)
       = f_\#P(Y)+g_\#P(Z).
 \tag{PI}\label{eq:akr-two-image-saturation}
\]
The same assertion holds with coefficients in an arbitrary abelian
group. Thus a protected-cycle combination whose edge support lies in
the union uses only protected cells factoring through one of those
two arrows.
\end{proposition}

\begin{proof}
The inclusion from right to left follows from functoriality. For the
other inclusion, use the unique expansion in the protected-cell
basis, and consider the nonempty family, if any, of its cells which
factor through neither $f$ nor $g$. Index those cells by
$i=1,\ldots,r$, write their marked images as $S_i$, and let
$s_i=|S_i|\geq45$. Each meets either full arrow image in at most two
points, by the protected routing theorem. At a stage $m$ containing
all the data, put $q=a_m/a_1>3/4$. The weighted counting inequality
gives
\[
 \left|\bigcup_i S_i\setminus f(Y)\right|
       \geq q\sum_i(s_i-2).
\]
Deleting the second image removes at most $2r$ further points.
Consequently
\[
 \left|\bigcup_i S_i\setminus\bigl(f(Y)\cup g(Z)\bigr)\right|
 \geq q\sum_i(s_i-2)-2r
 >\frac34\sum_i s_i-\frac72r
 >\frac12\sum_i s_i.
\]
The final strict inequality uses $s_i\geq45$. If every point of this
outside union belonged to at least two of the indicated marked
circles, the number of marked incidences would be at least twice
that union's size, contradicting the bound by $\sum_i s_i$. There is
therefore a marked vertex outside both images belonging to exactly
one of the indicated cells. All other cells in the full expansion
factor through $f$ or $g$ and have every marked vertex inside the
union, so this vertex belongs to no other cell in the expansion.

Every boundary vertex of a protected cell is marked. Either boundary
edge incident to this private vertex consequently belongs to that
cell alone, and lies outside $U$. Its coefficient in the given
combination is the nonzero cell coefficient, with sign $\pm1$.
This contradicts the assumed edge support in $U$. Thus no such cell
occurs, proving \eqref{eq:akr-two-image-saturation}. The argument only
uses that changing sign preserves a nonzero coefficient, so it works
for coefficients in any abelian group.
\end{proof}

\paragraph{Limitation of this local conclusion.}
Proposition~\ref{C:protected-two-images} is a statement about chains
at one object. It does not assert that the protected subgroup remains
injective after taking the colimit of $K$. A path-difference class
from Proposition~\ref{C:fixed-path-kernel} could agree with a protected
combination only after old colimit relations have been used, rather
than as chains at its target. Controlling those prior relations is
still required for the all-stage assertion in C.1.
